\subsection{Prior work}

Oneto and Ventura ask for the partially symmetric rank of
$W_{d_1}\otimes\cdots\otimes W_{d_k}$ with $d_i\geq3$ in
\cite[Question~12]{OV}. Here $W_d$ corresponds, up to a nonzero scalar,
to $x^{d-1}y$, and each copy remains a separate symmetric block.
Ballico, Bernardi, Christandl and Gesmundo proved
\[
 \Rk(W_3^{\otimes k})\leq(k+2)2^{k-1}
 \quad\text{and}\quad
 \Rk(W_{d_1}\otimes\cdots\otimes W_{d_k})
 \leq 2^{k-1}\sum_{i=1}^k d_i,
\]
the latter for $k\geq2$ and $d_i\geq3$
\cite[Theorems~3.3 and~3.6]{BBCG}. Their Proposition~4.3 gives the
lower bound $\sum_i d_i-k+1$ already for ordinary tensor rank, hence
also for partially symmetric rank. Canino, Casarotti and Santarsiero
subsequently obtained
\[
 \Rk(W_{d_1}\otimes\cdots\otimes W_{d_k})
 \leq 2^{k-1}\left(\sum_{i=1}^k d_i-2k+2\right)
 \qquad(k\geq2,\ d_i\geq3)
\]
by an explicit construction \cite[Theorem~1.1]{CCS}.

For two binary-monomial blocks, Ga{\l}\k{a}zka constructed a reduced
apolar set giving
\[
 \Rk(x_1^{a_1}y_1^{b_1}x_2^{a_2}y_2^{b_2})
 \leq (a_1+1)(b_2+1)+(b_1+1)(a_2+1)
          -(b_1+1)(b_2+1),
 \qquad a_i\geq b_i\geq1,
\]
and obtained lower bounds and exact special cases
\cite[Theorem~1.5 and Section~7.1 of the arXiv version]{Galazka}.
The binomial equations in that construction describe a finite set in
the two-dimensional torus. The lower bound
\[
 \Rk(x_1^{a_1}y_1^{b_1}x_2^{a_2}y_2^{b_2})
 \geq \max\{(a_1+1)(b_2+1),(b_1+1)(a_2+1)\}
\]
also appears in \cite[Proposition~4.3 of the arXiv version]{BBGV}.

Behr's unpublished manuscript \cite{Behr} gives the exact two-block
formula
\[
 \Rk(x_1^{a_1}y_1^{b_1}x_2^{a_2}y_2^{b_2})
 =(a_1+1)(a_2+1)-(a_1-b_1)(a_2-b_2)
 \qquad(a_i\geq b_i\geq1)
\]
in Corollary~5.4, including
$\Rk(W_{d_1}\otimes W_{d_2})=2(d_1+d_2-2)$ for $d_1,d_2\geq2$
in Theorem~1.1. Thus the two-block specialization is already present
in that work. Its section on three or more factors leaves that case open.
